This report details the implementation and results of the Advanced Analytics in Business course projects. The work covers four primary tasks: customer lifetime value prediction, image-based classification of Lego minifigures, the development of a Retrieval-Augmented Generation (RAG) system for Steam game recommendations, and multilayer network analysis of the Jeffrey Epstein document corpus.

\section{Code availability}
All source code, data preprocessing scripts, and model training notebooks are provided as supplementary materials in this submission. For convenience and version-controlled auditing, the complete repository is also hosted on GitHub at:
\begin{center}
    \url{https://github.com/BigDataNoSleep/AdvancedAnalytics}
\end{center}

\section{Project overview}
The report is structured into four main chapters, each corresponding to one of the primary assignments:
\begin{itemize}
    \item \textbf{Task 1: Customer lifetime value prediction.} This task involves predicting future customer revenue using historical transaction data, employing various exploratory data analysis techniques, feature engineering, and tree-based gradient boosting models.
    \item \textbf{Task 2: Lego minifigure theme classification.} This task focuses on computer vision, comparing MobileNetV2 and EfficientNetV2S backbones to classify images of Lego minifigures into specific themes.
    \item \textbf{Task 3: LLM-based recommender system.} This task presents a hybrid Retrieval-Augmented Generation system for Steam games, combining dense vector search, sparse BM25 retrieval, and a cross-encoder for reranking, with an interactive front-end.
    \item \textbf{Task 4: Epstein corpus network analysis.} This task applies multilayer network analysis to the Jeffrey Epstein document corpus, exploring social co-travel, bipartite person-location relationships, and financial transactions to map institutional and personal influence.
\end{itemize}
